\documentclass[11pt]{article}

\usepackage[margin=1in]{geometry}
\usepackage{amsmath, amssymb}
\usepackage{enumitem}
\usepackage{setspace}

\setlength{\parindent}{0pt}
\setlength{\parskip}{8pt}

\begin{document}

\begin{center}
    {\Large \textbf{Time Value of Money Problems}}\\
\end{center}

\vspace{10pt}

\textbf{Problem 1.} Consider the following amount function, which is a function of the number of years elapsed, $t$:
\[
A(t) = Xt^{2} + Yt + Z.
\]

You are given that:
\begin{itemize}[leftmargin=2em]
    \item $A(0) = 500$
    \item $A(1) = 521$
    \item $A(4) = 596$
\end{itemize}

Calculate the future value at time $5$ of \(\$600\) that is deposited at time \(0\).

\medskip

\hrule

\medskip


\textbf{Problem 2.} Fund A accumulates at a force of interest,
\[
\delta_t^{A} = \frac{0.05}{1 + 0.05t}, \qquad t \ge 0.
\]

Fund B accumulates at an annual force of interest of $0.05$.

You are given that:
\begin{itemize}[leftmargin=2em]
    \item The amount in Fund A at time $0$ is $1{,}000$.
    \item The amount in Fund B at time $0$ is $500$.
    \item The amount in Fund C at any time $t$ is equal to the sum of the amount in Fund A and the amount in Fund B.
    \item Fund C accumulates at a force of interest, $\delta_t^{C}$, $t \ge 0$.
\end{itemize}

Calculate $\delta_2^{C}$.

\medskip

\hrule

\medskip


\textbf{Problem 3.} Clayton has two savings accounts. Savings account \#1 earns interest at an annual effective discount rate of $4.7619\%$, and savings account \#2 earns interest at a nominal annual interest rate compounded monthly of $11.3861\%$. Clayton has not made any deposits or withdrawals since March 1, 2014, when the amount in savings account \#1 was five times the amount in savings account \#2.

The sum of the two accounts on March 1, 2025 was \$130{,}000.

Calculate the sum of the two accounts on March 1, 2014.

\medskip



\newpage


\textbf{Problem 4.} Three deposits are made into a fund as follows:
\begin{itemize}[leftmargin=2em]
    \item A deposit of \$100 now
    \item A deposit of \$300 in one year
    \item A deposit of \$500 in two years
\end{itemize}

An annual simple interest of $6\%$ is credited in the first and second year. For the third year, interest is credited at
\[
\delta_t = 0.06t,
\]
where $t$ is the number of years from now.

Find the accumulated value of the fund at the end of three years.

\medskip

\hrule

\medskip

\textbf{Problem 5.} The present value of \$200 paid at the end of $n$ years, plus the present value of \$100 paid at the end of $2n$ years, is \$200.

Determine the annual effective rate of interest.

\medskip

\begin{enumerate}[label=\textbf{\Alph*.}, leftmargin=2.5em]
    \item $\left(\dfrac{\sqrt{3}+1}{2}\right)^{1/n} - 1$
    
    \item $1 - \left(\dfrac{\sqrt{3}-1}{2}\right)^{1/n}$
    
    \item $\left(\dfrac{\sqrt{3}-1}{2}\right)^{1/n} - 1$
    
    \item $\left(\dfrac{\sqrt{3}+1}{2}\right)^{1/(2n)} - 1$
    
    \item $\left(\dfrac{\sqrt{3}-1}{2}\right)^{1/n} + 1$2
\end{enumerate}

\medskip

\hrule

\medskip

\newpage

\textbf{Problem A.} \$1{,}000 is deposited into a savings account at time $t=0$. No other amounts are deposited. The accumulated amount in the account at any time is given by
\[
A(t) = 1000\left(1 + \frac{2t}{35}\right)^2.
\]

At what time $t$ is the force of interest equal to $0.04$?

\medskip

\hrule

\medskip

\textbf{Problem B.} Joe deposits \$10 today and another \$30 in five years into a fund paying simple interest of $11\%$ per year.

Tina will make the same two deposits, but the \$10 will be deposited $n$ years from today and the \$30 will be deposited $2n$ years from today. Tina's deposits earn an annual effective interest rate of $9.15\%$.

At the end of $10$ years, the accumulated amount of Tina's deposits equals the accumulated amount of Joe's deposits.

Calculate $n$.

\medskip

\hrule


\medskip


\textbf{Problem C.} Jim deposits \$100 into an account that credits:
\begin{itemize}[leftmargin=2em]
    \item An effective rate of discount of $3\%$ for the first year;
    \item An effective rate of discount of $6\%$ for the second year;
    \item An effective rate of discount of $9\%$ for the third year; and
    \item $\delta_t = 0.02t$ for the fourth and fifth years.
\end{itemize}

How much is in the account at the end of the fifth year?

\medskip

\hrule





\end{document}
